\documentclass[12pt]{article}

%----------------- Packages ------------------
\usepackage[utf8]{inputenc}
\usepackage[T1]{fontenc}
\usepackage{amsmath, amssymb}
\usepackage{geometry}
\usepackage{xcolor}
\usepackage{titlesec}
\usepackage{fancyhdr}
\usepackage[breakable]{tcolorbox}
\usepackage{graphicx}
\usepackage{tikz}
\usepackage{float}
\usetikzlibrary{arrows.meta, decorations.markings}

%----------------- Page Setup -----------------
\geometry{a4paper, margin=2.5cm}
\pagestyle{fancy}
\fancyhf{}
\rhead{Final Exam — \textbf{2022}}
\lhead{Algebra IV}
\cfoot{\thepage}

%----------------- Title Styling --------------
\titleformat{\section}{\normalfont\Large\bfseries}{Exercise \thesection:}{1em}{}
\titleformat{\subsection}{\normalfont\bfseries}{Answer:}{1em}{}

%----------------- Custom Boxes ----------------
\tcbuselibrary{listingsutf8}
\newtcolorbox{answerbox}{
  colback=gray!10,
  colframe=black,
  fonttitle=\bfseries,
  title=Answer Area,
  breakable,
  before skip=10pt,
  after skip=10pt
}

%----------------- Document Start --------------
\begin{document}

%----------------- Exam Info -------------------
\begin{center}
  \Large\textbf{Abdelmalek Essadi University} \\[1em]
  \large\textit{Algebra IV — Final Exam} \\[0.5em]
  \large\textit{Year: 2022} \\[2em]
\end{center}

\vspace{0.5cm}

%----------------- EXERCISE 1 ------------------
\section{}
% student's answer area
Let $\alpha\in\mathbb{R}$. Consider the matrix $A_{\alpha}\in M_{3}(\mathbb{R})$
defined by: 
$$A_{\alpha} = 
\begin{pmatrix} 
    1 & 1 & 1 \\ 
    1 & 1 & 1 \\ 
    -\alpha & 1 & 1 
\end{pmatrix}$$ 

\begin{enumerate} 
    \item Determine the \textbf{characteristic polynomial} of the matrix $A_{\alpha}$. 
    \item For which values of $\alpha$ is the matrix $A_{\alpha}$ \textbf{diagonalizable} over $\mathbb{R}$? 
    \item Consider the equation
in $M_3(\mathbb{R})$: $$M^2 = A_{\alpha} \quad (*)$$ 
For which values of $\alpha$ does this equation $(*)$ possess \textbf{4 diagonalizable
solutions} over $\mathbb{R}$? 
\end{enumerate}

\begin{answerbox}

\end{answerbox}

%----------------- EXERCISE 2 ------------------
\section{}

Consider the matrix $A\in M_{3}(\mathbb{R})$, defined by: 
$$A = 
\begin{pmatrix}
-1 & -2 & 0 \\ 
2 & 3 & 0 \\ 
2 & 2 & 1 
\end{pmatrix}$$ 

\begin{enumerate}
\item Determine a \textbf{Jordan reduction relation} for the matrix $A$. 
\item Deduce a \textbf{Jordan reduction relation} for the matrix $-A^{-1}$. 
\item Determine the polynomials $\mathcal{Q}\in\mathbb{R}[X]$ such that the matrix $\mathcal{Q}(-A^{-1})$ is \textbf{diagonalizable}. 
\end{enumerate}

\begin{answerbox}

\end{answerbox}

%----------------- EXERCISE 3 ------------------
\section{}
Let $\alpha \in \mathbb{R}$. Consider the matrix $A_{\alpha} \in M_{3}(\mathbb{R})$, defined by:
$$A_{\alpha} = \begin{pmatrix} 2\alpha + 2 & \alpha + 1 & 1 \\ -4\alpha - 2 & -2\alpha - 1 & -2 \\ 2\alpha - 1 & \alpha - 1 & 0 \end{pmatrix}$$

\vspace{0.5cm}

\begin{enumerate}
    \item \begin{enumerate}
        \item Determine $\mathrm{Rg}(A_{\alpha} - I_3)$.
        \item Deduce $N_{\alpha}(X)$.
        \item Without performing any calculation, what is the value of $\dim E_{-1}(A_{\alpha})$?
    \end{enumerate}
    
    \item For which values of $\alpha \in \mathbb{R}$ is the matrix $A_{\alpha}$ \textbf{diagonalizable}?
    
    \item Write the \textbf{minimal polynomial} of $A_{\alpha}$.
    
    \item We choose $\alpha = 1$ and set $A = A_1$.
    \begin{enumerate}
        \item Determine a \textbf{Jordan reduction relation} for $A$.
        \item Determine the polynomials $P \in \mathbb{R}[X]$ such that $A$ and $P(A)$ are \textbf{similar}.
    \end{enumerate}
    
    \item \begin{enumerate}
        \item Determine a \textbf{diagonalization relation} for $A_0$.
        \item Without performing any calculation, what is the value of $A_0^k$, where $k \in \mathbb{N}$?
        \item Determine the matrices $M \in M_{3}(\mathbb{R})$ that satisfy $M^2 = A_0 + I_3$ and $\mathrm{Tr}(M) > 0$.
    \end{enumerate}
\end{enumerate}

\begin{answerbox}

\end{answerbox}

%----------------- END ------------------
\end{document}